% ==============================================================================
% FICHIER : introduction_generale.tex
% Description : Introduction Générale du rapport SafeCity-Connect
% Auteurs : Dhia Fersi & Mohamed Iyed Amiri (ISET Bizerte)
% Organisme d'accueil : HRMAPS
% ==============================================================================

\chapter*{Introduction Générale}
\addcontentsline{toc}{chapter}{Introduction Générale}

La progression rapide du paysage technologique contemporain a propulsé la transformation numérique au rang d'enjeu stratégique majeur pour l'ensemble des organisations, qu'elles relèvent du secteur public ou privé. Adopter des solutions informatiques performantes constitue aujourd'hui un levier indispensable pour se défaire des processus archaïques, accélérer le traitement de l'information et offrir aux citoyens des prestations fiables, agiles et en adéquation avec leurs attentes croissantes.

À cet égard, la gouvernance des espaces urbains et la gestion du domaine public appellent la mise en place de dispositifs technologiques novateurs favorisant l'implication collective. Les autorités municipales se trouvent en effet confrontées, au quotidien, à un éventail de dysfonctionnements sur la voie publique~-- nids-de-poule, ruptures de canalisations, pannes du réseau d'éclairage ou accumulations illicites de déchets~-- qui altèrent significativement la qualité de vie des résidents tout en générant des charges d'entretien considérables. Le paradigme des villes intelligentes (Smart Cities) offre des solutions tangibles à ces défis en promouvant une approche participative de la gestion urbaine, où le citoyen devient un acteur central du signalement et où les responsables territoriaux disposent d'outils analytiques robustes pour piloter les interventions.

C'est précisément dans ce contexte que s'inscrit notre projet de fin d'études, réalisé au sein de l'entreprise \textbf{HRMAPS}. L'objectif principal a porté sur la conception et l'élaboration de \textbf{SafeCity-Connect} : une plateforme de signalement géolocalisé et de suivi des incidents urbains. Elle combine un \textbf{backend Spring Boot} monolithique, un portail \textbf{Angular~18} (carte publique sans connexion, espaces citoyen, administrateur et services techniques), une application \textbf{Flutter}, l'authentification \textbf{Keycloak}, une base \textbf{PostgreSQL/PostGIS} et un service \textbf{FastAPI/YOLOv8} pour la classification automatique des photos.

Ce rapport retrace de façon détaillée l'ensemble des étapes de réalisation et de maturation de ce projet. Sa structure se décline selon l'articulation suivante :

Le premier chapitre, \textbf{« Présentation du cadre du projet »}, expose l'organisme d'accueil \textbf{HRMAPS} ainsi que son fonctionnement interne. Il dresse un état des lieux critique des pratiques existantes afin de mettre en lumière leurs limites et de légitimer la solution envisagée. Ce chapitre précise en outre les orientations méthodologiques retenues (associant Scrum et UML), les ressources matérielles mobilisées et le schéma d'architecture n-tiers adopté.

Le deuxième chapitre, intitulé \textbf{« Analyse et spécification des besoins »}, procède au recensement des acteurs impliqués, à la formalisation des exigences fonctionnelles et non fonctionnelles du système, et à l'ordonnancement des travaux au travers du carnet de produit (Product Backlog).

Le troisième chapitre, dédié à la \textbf{« Release 1 : Gestion d'accès, profils et sécurité »}, décrit l'édification du socle sécuritaire centralisé fondé sur \textbf{Keycloak}, l'administration des profils utilisateurs, et les premières briques d'implémentation (Sprint 0 et livrables fondamentaux).

Le quatrième chapitre, portant sur la \textbf{« Release 2 : Gestion avancée, rapports et intelligence artificielle »}, traite de l'affectation aux services techniques, des preuves de réparation, des tableaux de bord et heatmap, de la classification YOLOv8, des alertes administrateur dédiées aux nouveaux reports à vérifier, du support par tickets et de la gamification (+10 points à la validation).

Pour clore ce mémoire, une conclusion générale synthétise les résultats obtenus au regard des objectifs fixés et esquisse les axes d'amélioration et d'extension envisageables pour le projet.
